\documentclass[10pt,a4paper]{amsart}
\usepackage{amsmath,amssymb,amsthm,graphicx}
% Retained minimal preamble for the Stacks Project chapter more-algebra.tex (More on Algebra).
% The chapter file pinned-original-source.tex carries only \input{preamble} on its first line, so no
% preamble text lives inside the chapter itself.  The Stacks Project uses one single shared preamble for
% every chapter; the snapshot stages it identically as sources/stacks-chow/preamble.tex and
% sources/stacks-etale-cohomology/preamble.tex.  The block below is copied mechanically, line for line,
% from sources/stacks-chow/preamble.tex lines 14-18 (the Xy-pic setup) and 154-217 (the native theorem
% declarations together with the mathematical macros).  Nothing was added, removed or redefined; the
% wrapper class and the amsmath/amssymb/amsthm/graphicx base packages are supplied by the assembler.
% Chapter: more-algebra.tex -- More on Algebra.
% For commutative diagrams we use Xy-pic
\usepackage[all]{xy}

% We use 2cell for 2-commutative diagrams.
\xyoption{2cell}

% Theorem environments.
%
\theoremstyle{plain}
\newtheorem{theorem}[subsection]{Theorem}
\newtheorem{proposition}[subsection]{Proposition}
\newtheorem{lemma}[subsection]{Lemma}

\theoremstyle{definition}
\newtheorem{definition}[subsection]{Definition}
\newtheorem{example}[subsection]{Example}
\newtheorem{exercise}[subsection]{Exercise}
\newtheorem{situation}[subsection]{Situation}

\theoremstyle{remark}
\newtheorem{remark}[subsection]{Remark}
\newtheorem{remarks}[subsection]{Remarks}

\numberwithin{equation}{subsection}

% Macros
%
\def\lim{\mathop{\mathrm{lim}}\nolimits}
\def\colim{\mathop{\mathrm{colim}}\nolimits}
\def\Spec{\mathop{\mathrm{Spec}}}
\def\Hom{\mathop{\mathrm{Hom}}\nolimits}
\def\Ext{\mathop{\mathrm{Ext}}\nolimits}
\def\SheafHom{\mathop{\mathcal{H}\!\mathit{om}}\nolimits}
\def\SheafExt{\mathop{\mathcal{E}\!\mathit{xt}}\nolimits}
\def\Sch{\mathit{Sch}}
\def\Mor{\mathop{\mathrm{Mor}}\nolimits}
\def\Ob{\mathop{\mathrm{Ob}}\nolimits}
\def\Sh{\mathop{\mathit{Sh}}\nolimits}
\def\NL{\mathop{N\!L}\nolimits}
\def\CH{\mathop{\mathrm{CH}}\nolimits}
\def\proetale{{pro\text{-}\acute{e}tale}}
\def\etale{{\acute{e}tale}}
\def\QCoh{\mathit{QCoh}}
\def\Ker{\mathop{\mathrm{Ker}}}
\def\Im{\mathop{\mathrm{Im}}}
\def\Coker{\mathop{\mathrm{Coker}}}
\def\Coim{\mathop{\mathrm{Coim}}}

% Boxtimes
%
\DeclareMathSymbol{\boxtimes}{\mathbin}{AMSa}{"02}

%
% Macros for moduli stacks/spaces
%
\def\QCohstack{\mathcal{QC}\!\mathit{oh}}
\def\Cohstack{\mathcal{C}\!\mathit{oh}}
\def\Spacesstack{\mathcal{S}\!\mathit{paces}}
\def\Quotfunctor{\mathrm{Quot}}
\def\Hilbfunctor{\mathrm{Hilb}}
\def\Curvesstack{\mathcal{C}\!\mathit{urves}}
\def\Polarizedstack{\mathcal{P}\!\mathit{olarized}}
\def\Complexesstack{\mathcal{C}\!\mathit{omplexes}}
% \Pic is the operator that assigns to X its picard group, usage \Pic(X)
% \Picardstack_{X/B} denotes the Picard stack of X over B
% \Picardfunctor_{X/B} denotes the Picard functor of X over B
\def\Pic{\mathop{\mathrm{Pic}}\nolimits}
\def\Picardstack{\mathcal{P}\!\mathit{ic}}
\def\Picardfunctor{\mathrm{Pic}}
\def\Deformationcategory{\mathcal{D}\!\mathit{ef}}
\usepackage{cleveref}
\crefname{originallemma}{original lemma}{original lemmas}
\Crefname{originallemma}{Original lemma}{Original lemmas}
\makeatletter
\newlabel{algebra-lemma-finite-over-complete-ring}{{031D}{source}{}{sourcepointer-1}{}}
\newlabel{algebra-lemma-finite-over-complete-ring@cref}{{[originallemma][0][]031D}{[1][0][]source}{}{}{}}
\makeatother
\crefname{originallemma}{original lemma}{original lemmas}
\Crefname{originallemma}{Original lemma}{Original lemmas}
\makeatletter
\newlabel{algebra-theorem-cohen-structure-theorem}{{032A}{source}{}{sourcepointer-2}{}}
\newlabel{algebra-theorem-cohen-structure-theorem@cref}{{[originallemma][0][]032A}{[1][0][]source}{}{}{}}
\makeatother
\crefname{originaldefinition}{original definition}{original definitions}
\Crefname{originaldefinition}{Original definition}{Original definitions}
\makeatletter
\newlabel{definition-solution}{{34445}{source}{}{sourcepointer-3}{}}
\newlabel{definition-solution@cref}{{[originaldefinition][0][]34445}{[1][0][]source}{}{}{}}
\makeatother
\crefname{originallemma}{original lemma}{original lemmas}
\Crefname{originallemma}{Original lemma}{Original lemmas}
\makeatletter
\newlabel{lemma-abhyankar}{{34037}{source}{}{sourcepointer-4}{}}
\newlabel{lemma-abhyankar@cref}{{[originallemma][0][]34037}{[1][0][]source}{}{}{}}
\makeatother
\crefname{originallemma}{original lemma}{original lemmas}
\Crefname{originallemma}{Original lemma}{Original lemmas}
\makeatletter
\newlabel{lemma-approximate-separable-extension}{{33814}{source}{}{sourcepointer-5}{}}
\newlabel{lemma-approximate-separable-extension@cref}{{[originallemma][0][]33814}{[1][0][]source}{}{}{}}
\makeatother
\crefname{originallemma}{original lemma}{original lemmas}
\Crefname{originallemma}{Original lemma}{Original lemmas}
\makeatletter
\newlabel{lemma-extension-dvrs-formally-smooth}{{32989}{source}{}{sourcepointer-6}{}}
\newlabel{lemma-extension-dvrs-formally-smooth@cref}{{[originallemma][0][]32989}{[1][0][]source}{}{}{}}
\makeatother
\crefname{originallemma}{original lemma}{original lemmas}
\Crefname{originallemma}{Original lemma}{Original lemmas}
\makeatletter
\newlabel{lemma-formally-smooth-goes-up}{{34006}{source}{}{sourcepointer-7}{}}
\newlabel{lemma-formally-smooth-goes-up@cref}{{[originallemma][0][]34006}{[1][0][]source}{}{}{}}
\makeatother
\crefname{originallemma}{original lemma}{original lemmas}
\Crefname{originallemma}{Original lemma}{Original lemmas}
\makeatletter
\newlabel{lemma-lift-continuous}{{9257}{source}{}{sourcepointer-8}{}}
\newlabel{lemma-lift-continuous@cref}{{[originallemma][0][]9257}{[1][0][]source}{}{}{}}
\makeatother
\crefname{originallemma}{original lemma}{original lemmas}
\Crefname{originallemma}{Original lemma}{Original lemmas}
\makeatletter
\newlabel{lemma-pull-root-uniformizer}{{33954}{source}{}{sourcepointer-9}{}}
\newlabel{lemma-pull-root-uniformizer@cref}{{[originallemma][0][]33954}{[1][0][]source}{}{}{}}
\makeatother
\crefname{originallemma}{original lemma}{original lemmas}
\Crefname{originallemma}{Original lemma}{Original lemmas}
\makeatletter
\newlabel{lemma-solution-after-strict-henselization}{{34555}{source}{}{sourcepointer-10}{}}
\newlabel{lemma-solution-after-strict-henselization@cref}{{[originallemma][0][]34555}{[1][0][]source}{}{}{}}
\makeatother
\crefname{originallemma}{original lemma}{original lemmas}
\Crefname{originallemma}{Original lemma}{Original lemmas}
\makeatletter
\newlabel{lemma-solutions-go-down}{{34517}{source}{}{sourcepointer-11}{}}
\newlabel{lemma-solutions-go-down@cref}{{[originallemma][0][]34517}{[1][0][]source}{}{}{}}
\makeatother
\crefname{originallemma}{original lemma}{original lemmas}
\Crefname{originallemma}{Original lemma}{Original lemmas}
\makeatletter
\newlabel{lemma-special-case}{{35370}{source}{}{sourcepointer-12}{}}
\newlabel{lemma-special-case@cref}{{[originallemma][0][]35370}{[1][0][]source}{}{}{}}
\makeatother
\crefname{originalremark}{original remark}{original remarks}
\Crefname{originalremark}{Original remark}{Original remarks}
\makeatletter
\newlabel{remark-construction}{{33882}{source}{}{sourcepointer-13}{}}
\newlabel{remark-construction@cref}{{[originalremark][0][]33882}{[1][0][]source}{}{}{}}
\makeatother
\title{Epp's theorem: a finite weak solution exists for an extension of discrete valuation rings when the p-power core of the bigger residue field is separable over the smaller}
\author{The Stacks Project Authors}
\date{stacks-project:more-algebra.tex:local-snapshot-20261012; GNU FDL 1.2 or later}
\begin{document}
\maketitle
\setcounter{subsection}{1}
\noindent\textit{Editorial note.} All delivered mathematical text is the author's own, copied from the pinned Stacks Project chapter file more-algebra.tex; no statement or proof step was rewritten, reordered or removed.\par\medskip
\noindent\textit{Editorial note.} The chapter file pinned-original-source.tex carries only \textbackslash{}input\{preamble\} on its first line, so no preamble text lives inside the chapter. The retained file sources/stacks-more-algebra/retained-preamble.tex is a mechanical, line-for-line copy of the project's single shared preamble as staged in the snapshot; nothing was added, removed or redefined and no package was added.\par\medskip
\noindent\textit{Editorial note.} The source URL is the source-level Stacks Project URL, as required when more than one card is produced from one source; the per-claim tag is carried by claim\_key.\par\medskip
\noindent\textit{Editorial note.} The references inside the counted proof span are recorded as non-load-bearing navigation pointers; they name the case lemmas and standard structural facts whose applications together with the surrounding argument stay inside the counted proof span.\par\medskip
\noindent\textit{Allowed standard prerequisites (editorial).}\begin{itemize}
\item Abhyankar's lemma on ramification and the construction of extensions adjoining a root of a uniformizer. invoked at proof lines 35533-35543. Settles the tame case and supplies the formal smoothness conclusion used there.
\item Strict henselization of a discrete valuation ring and its effect on the residue field. invoked at proof lines 35547-35555. Is the first reduction, replacing the residue fields by their separable closures.
\item Completions of discrete valuation rings, formal smoothness and the descent of solutions. used at proof lines 35558-35586. Is the second reduction, moving the problem to the complete case and then descending the answer.
\item The Cohen structure theorem and Cohen rings. invoked at proof lines 35600-35605. Produces the ring Lambda mapping to A with t sent to pi, which the explicit construction in the complete case uses.
\item Weak solutions for an extension of discrete valuation rings. the definition invoked at statement lines 35526-35527. Fixes the property that K\_1/K is required to have.
\end{itemize}\medskip
\ifdefined\hypertarget\hypertarget{sourcepointer-1}{}\fi
\noindent\textit{Source reference.} Original lemma 031D: focused-seeds/AGP/sources/stacks-snapshot-20261011/tags.txt, source lines 3980--3980. Cross-chapter navigation pointer into the Stacks Project; the standard fact it supplies is applied inside the counted proof span and the argument resting on it stays there.\par\medskip
\ifdefined\hypertarget\hypertarget{sourcepointer-2}{}\fi
\noindent\textit{Source reference.} Original lemma 032A: focused-seeds/AGP/sources/stacks-snapshot-20261011/tags.txt, source lines 4013--4013. Cross-chapter navigation pointer into the Stacks Project; the standard fact it supplies is applied inside the counted proof span and the argument resting on it stays there.\par\medskip
\ifdefined\hypertarget\hypertarget{sourcepointer-3}{}\fi
\noindent\textit{Source reference.} Original definition 34445: sources/stacks-more-algebra/pinned-original-source.tex, source lines 34445--34445. Same-chapter navigation pointer; the statement it names is standard for the section and the argument that consumes it stays inside the counted proof span.\par\medskip
\ifdefined\hypertarget\hypertarget{sourcepointer-4}{}\fi
\noindent\textit{Source reference.} Original lemma 34037: sources/stacks-more-algebra/pinned-original-source.tex, source lines 34037--34037. Same-chapter navigation pointer; the statement it names is standard for the section and the argument that consumes it stays inside the counted proof span.\par\medskip
\ifdefined\hypertarget\hypertarget{sourcepointer-5}{}\fi
\noindent\textit{Source reference.} Original lemma 33814: sources/stacks-more-algebra/pinned-original-source.tex, source lines 33814--33814. Same-chapter navigation pointer; the statement it names is standard for the section and the argument that consumes it stays inside the counted proof span.\par\medskip
\ifdefined\hypertarget\hypertarget{sourcepointer-6}{}\fi
\noindent\textit{Source reference.} Original lemma 32989: sources/stacks-more-algebra/pinned-original-source.tex, source lines 32989--32989. Same-chapter navigation pointer; the statement it names is standard for the section and the argument that consumes it stays inside the counted proof span.\par\medskip
\ifdefined\hypertarget\hypertarget{sourcepointer-7}{}\fi
\noindent\textit{Source reference.} Original lemma 34006: sources/stacks-more-algebra/pinned-original-source.tex, source lines 34006--34006. Same-chapter navigation pointer; the statement it names is standard for the section and the argument that consumes it stays inside the counted proof span.\par\medskip
\ifdefined\hypertarget\hypertarget{sourcepointer-8}{}\fi
\noindent\textit{Source reference.} Original lemma 9257: sources/stacks-more-algebra/pinned-original-source.tex, source lines 9257--9257. Same-chapter navigation pointer; the statement it names is standard for the section and the argument that consumes it stays inside the counted proof span.\par\medskip
\ifdefined\hypertarget\hypertarget{sourcepointer-9}{}\fi
\noindent\textit{Source reference.} Original lemma 33954: sources/stacks-more-algebra/pinned-original-source.tex, source lines 33954--33954. Same-chapter navigation pointer; the statement it names is standard for the section and the argument that consumes it stays inside the counted proof span.\par\medskip
\ifdefined\hypertarget\hypertarget{sourcepointer-10}{}\fi
\noindent\textit{Source reference.} Original lemma 34555: sources/stacks-more-algebra/pinned-original-source.tex, source lines 34555--34555. Same-chapter navigation pointer; the statement it names is standard for the section and the argument that consumes it stays inside the counted proof span.\par\medskip
\ifdefined\hypertarget\hypertarget{sourcepointer-11}{}\fi
\noindent\textit{Source reference.} Original lemma 34517: sources/stacks-more-algebra/pinned-original-source.tex, source lines 34517--34517. Same-chapter navigation pointer; the statement it names is standard for the section and the argument that consumes it stays inside the counted proof span.\par\medskip
\ifdefined\hypertarget\hypertarget{sourcepointer-12}{}\fi
\noindent\textit{Source reference.} Original lemma 35370: sources/stacks-more-algebra/pinned-original-source.tex, source lines 35370--35370. Same-chapter navigation pointer; the statement it names is standard for the section and the argument that consumes it stays inside the counted proof span.\par\medskip
\ifdefined\hypertarget\hypertarget{sourcepointer-13}{}\fi
\noindent\textit{Source reference.} Original remark 33882: sources/stacks-more-algebra/pinned-original-source.tex, source lines 33882--33882. Same-chapter navigation pointer; the statement it names is standard for the section and the argument that consumes it stays inside the counted proof span.\par\medskip
\begin{theorem}[Epp]
\label{theorem-epp}
Let $A \subset B$ be an extension of discrete valuation rings with fraction
fields $K \subset L$. If the characteristic of $\kappa_A$ is $p > 0$,
assume that every element of
$$
\bigcap\nolimits_{n \geq 1} \kappa_B^{p^n}
$$
is separable algebraic over $\kappa_A$. Then there exists a finite extension
$K_1/K$ which is a weak solution for $A \to B$ as defined in
Definition \ref{definition-solution}.
\end{theorem}

\begin{proof}
If the characteristic of $\kappa_A$ is zero or if the residue characteristic
is $p$, the ramification index is prime to $p$, and the residue field
extension is separable, then this follows from Abhyankar's lemma
(Lemma \ref{lemma-abhyankar}). Namely, suppose the ramification index
is $e$. Choose a uniformizer $\pi \in A$. Let $K_1/K$
be the extension obtained by adjoining an $e$th root of $\pi$.
By Lemma \ref{lemma-pull-root-uniformizer} we see that the integral
closure $A_1$ of $A$ in $K_1$ is a discrete valuation ring with
ramification index over $A$. Thus $A_1 \to (B_1)_\mathfrak m$
is formally smooth in the $\mathfrak m$-adic topology
for all maximal ideals $\mathfrak m$ of $B_1$
by Lemma \ref{lemma-abhyankar} and a fortiori these are weakly
unramified extensions of discrete valuation rings.

\medskip\noindent
From now on we let $p$ be a prime number and we assume that $\kappa_A$ has
characteristic $p$. We first apply
Lemma \ref{lemma-solution-after-strict-henselization}
to reduce to the case that $A$ and $B$ have separably closed residue fields.
Since $\kappa_A$ and $\kappa_B$ are replaced by their separable algebraic
closures by this procedure we see that we obtain
$$
\kappa_A \supset \bigcap\nolimits_{n \geq 1} \kappa_B^{p^n}
$$
from the condition of the theorem.

\medskip\noindent
Let $\pi \in A$ be a uniformizer. Let $A^\wedge$ and $B^\wedge$ be the
completions of $A$ and $B$. We have a commutative diagram
$$
\xymatrix{
B \ar[r] & B^\wedge \\
A \ar[u] \ar[r] & A^\wedge \ar[u]
}
$$
of extensions of discrete valuation rings. Let $K^\wedge$ be the fraction
field of $A^\wedge$. Suppose that we can find a finite extension
$M/K^\wedge$ which is (a) a weak solution for $A^\wedge \to B^\wedge$
and (b) a compositum of a separable extension and an extension obtained
by adjoining a $p$-power root of $\pi$. Then by
Lemma \ref{lemma-approximate-separable-extension}
we can find a finite extension $K_1/K$ such that
$K^\wedge \otimes_K K_1 = M$. Let $A_1$, resp.\ $A_1^\wedge$
be the integral closure of $A$, resp.\ $A^\wedge$ in $K_1$, resp.\ $M$.
Since $A \to A^\wedge$ is formally smooth in the $\mathfrak m^\wedge$-adic
topology
(Lemma \ref{lemma-extension-dvrs-formally-smooth})
we see that $A_1 \to A_1^\wedge$ is formally smooth in the
$\mathfrak m_1^\wedge$-adic topology
(Lemma \ref{lemma-formally-smooth-goes-up} and
$A_1$ and $A_1^\wedge$ are discrete valuation
rings by discussion in Remark \ref{remark-construction}).
We conclude from Lemma \ref{lemma-solutions-go-down} part (2)
that $K_1/K$ is a weak solution for $A \to B^\wedge$.
Applying Lemma \ref{lemma-solutions-go-down} part (1)
we see that $K_1/K$ is a weak solution for $A \to B$.

\medskip\noindent
Thus we may assume $A$ and $B$ are complete discrete valuation rings
with separably closed residue fields of characteristic $p$ and
with $\kappa_A \supset \bigcap\nolimits_{n \geq 1} \kappa_B^{p^n}$.
We are also given a uniformizer $\pi \in A$ and we have to find a
weak solution for $A \to B$ which is a compositum of a separable
extension and a field obtained by taking $p$-power roots of $\pi$.
Note that the second condition is automatic if $A$ has mixed characteristic.

\medskip\noindent
Set $k = \bigcap\nolimits_{n \geq 1} \kappa_B^{p^n}$.
Observe that $k$ is an algebraically closed field of characteristic $p$.
If $A$ has mixed characteristic let $\Lambda$ be a Cohen
ring for $k$ and in the equicharacteristic case set $\Lambda = k[[t]]$.
We can choose a ring map $\Lambda \to A$ which maps $t$ to $\pi$ in the
equicharacteristic case. In the equicharacteristic case this follows
from the Cohen structure theorem
(Algebra, Theorem \ref{algebra-theorem-cohen-structure-theorem}) and
in the mixed characteristic case this follows as
$\mathbf{Z}_p \to \Lambda$ is formally smooth in the adic topology
(Lemmas \ref{lemma-extension-dvrs-formally-smooth} and
\ref{lemma-lift-continuous}).
Applying Lemma \ref{lemma-solutions-go-down} we see that it suffices to prove
the existence of a weak solution for $\Lambda \to B$ which in
the equicharacteristic $p$ case is a compositum of a separable
extension and a field obtained by taking $p$-power roots of $t$.
However, since $\Lambda = k[[t]]$ in the equicharacteristic case
and any extension of $k((t))$ is such a compositum, we can now
drop this requirement!

\medskip\noindent
Thus we arrive at the situation where $A$ and $B$ are complete, the residue
field $k$ of $A$ is algebraically closed of characteristic $p > 0$,
we have $k = \bigcap \kappa_B^{p^n}$, and in the mixed characteristic
case $p$ is a uniformizer of $A$ (i.e., $A$ is a Cohen ring for $k$).
If $A$ has mixed characteristic choose a Cohen ring
$\Lambda$ for $\kappa_B$ and in the equicharacteristic case set
$\Lambda = \kappa_B[[t]]$. Arguing as above we may choose a ring map
$A \to \Lambda$ lifting $k \to \kappa_B$ and mapping a uniformizer
to a uniformizer. Since $k \subset \kappa_B$ is separable the ring
map $A \to \Lambda$ is formally smooth in the adic topology
(Lemma \ref{lemma-extension-dvrs-formally-smooth}). Hence
we can find a ring map $\Lambda \to B$ such that the composition
$A \to \Lambda \to B$ is the given ring map $A \to B$ (see
Lemma \ref{lemma-lift-continuous}).
Since $\Lambda$ and $B$ are complete discrete valuation rings with the same
residue field, $B$ is finite over $\Lambda$
(Algebra, Lemma \ref{algebra-lemma-finite-over-complete-ring}).
This reduces us to the special case discussed in
Lemma \ref{lemma-special-case}.
\end{proof}

\end{document}
